%% notes-tex/publishing-system/sections/2-manuscript.tex

\section{The manuscript build\secstatus{todo}}
\label{sec:manuscript}

%% SOURCE: notes/publishing-system.mermaid.paper-build.md, rendered by
%% notes/assets/publish/scripts/mermaid_pdf.sh and scaled to the text block by
%% `make diagrams`; the targets and outputs are those of paper/nature-biotech/Makefile.
\tcfig{figures/paper-build.pdf}{What \texttt{make} builds for the manuscript.
Yellow boxes are source files, whether edited by hand or written by a script;
orange boxes are make targets; purple boxes are built PDFs; blue boxes are outside
the repository; red boxes are publish steps. One shared body reaches three views
through three wrappers, and only the drafting view goes to Zotero.}{fig:paper}

\notesec{One body, three views}
\texttt{content.tex} holds the manuscript. \texttt{submission.tex},
\texttt{editing.tex} and \texttt{twocolumn.tex} each include it and differ only in
class options and in what the editing-only macros expand to: \texttt{\textbackslash
secstatus} draws a chip in the editing view and nothing in the other two, and the
word-budget tags are counted only there. So \texttt{make paper} produces three PDFs
from one set of edits, and \texttt{make editing} alone is the fast loop while
writing.

%% SOURCE: hand-authored from paper/nature-biotech/Makefile.
\begin{table}[htbp]
\centering
\footnotesize
\caption{Manuscript targets. Every view depends on \texttt{checkfigs}, so a figure
out of spec stops all three builds before Tectonic runs.}
\label{tab:papertargets}
\begin{tabular}{lp{0.62\textwidth}}
\toprule
Target & Produces \\
\midrule
\texttt{make paper} & \texttt{submission.pdf}, \texttt{editing.pdf}, \texttt{twocolumn.pdf}, then the word-count report \\
\texttt{make submission}, \texttt{editing}, \texttt{twocolumn} & one view \\
\texttt{make figproto} & \texttt{figure-proto.pdf}, the true-scale figure-sizing canvas \\
\texttt{make figlimits} & \texttt{figure-limits.pdf}, the Nature print-box card for collaborators \\
\texttt{make all} & paper, figproto and figlimits \\
\texttt{make figures} & re-exports every referenced figure whose draw.io source is newer than its PDF \\
\texttt{make fig} & re-exports every referenced figure regardless of timestamps \\
\texttt{make checkfigs} & the size and scale gate on \texttt{figures/*.pdf}, exit non-zero on a violation \\
\texttt{make wordcount} & \texttt{wordcount.tex}, the per-section counts against the budgets \\
\texttt{make bib-pull} & \texttt{references.bib} from the served \texttt{paper} bibliography \\
\texttt{make publish} & \texttt{editing.pdf} into Zotero; \texttt{publish-dry}, \texttt{publish-list}, \texttt{publish-submission} \\
\texttt{make flat} & \texttt{submission-flat.tex}, one file for the journal's upload \\
\bottomrule
\end{tabular}
\end{table}

\notesec{Figures}
Every manuscript figure is composed in draw.io at Nature print size and kept as an
editable SVG under \texttt{notes/assets/drawio/NAME.drawio.svg}. \texttt{make
figures} exports it headless to \texttt{figures/{}NAME.pdf}, the same basename, so
source and export are matched by name and nothing has to be listed. The set of
figures to export is discovered from the \texttt{.tex}: any \texttt{figures/{}NAME.pdf}
the body references that has a draw.io source on disk. A figure whose MathJax
label rendered wider than its cell gets a crop-to-ink pass so the page lands at its
true ink box without scaling the fonts. \texttt{check-figures.sh} then measures
every PDF against the 180 by 240 mm print box and refuses any figure placed with a
scaling macro; because every view depends on it, an out-of-spec figure cannot reach
a PDF.

\notesec{Word budgets}
\texttt{wordcount.py} strips figure and table environments and the budget tags,
counts the rest with \texttt{texcount}, and writes \texttt{wordcount.tex}, which the
editing view embeds on its last page. \texttt{editing.pdf} depends on that file, so
the embedded counts always match the compiled draft.

\notesec{Two copies}
The repository directory is the workshop: private, full, with the drafting
scaffolding. \texttt{sync-overleaf.sh} copies a curated list of files into the
Overleaf-backed clone that collaborators see, with \texttt{submission.tex}
published as \texttt{main.tex}. It pulls first, so collaborators' additions
survive, and it propagates the workshop's own deletions through a manifest without
touching files it did not publish.
